\chapter{Kidney Infection (Pyelonephritis)}
\index{Kidney Infection (Pyelonephritis)}

\section*{Definition and Overview}
A kidney infection, medically called pyelonephritis, is an infection of one or both kidneys, usually caused by bacteria that have travelled up from the bladder through the ureters. It is more serious than a simple bladder infection (cystitis) and causes fever, chills, flank pain, and a feeling of being unwell. Kidney infections are most common in women, and risk is increased by pregnancy, kidney stones, urinary obstruction, catheters, and conditions that affect the flow of urine. Treatment is with antibiotics, often starting in hospital, and most people recover fully, though complications such as abscess and sepsis can occur. Recurrent infections and chronic kidney damage are reduced by treating underlying causes.

\section*{Epidemiology}
Pyelonephritis affects around 1 to 2 per cent of women each year and is much less common in men. It accounts for a significant number of hospital admissions, and it is one of the most common serious infections in pregnancy, where it can affect both mother and baby. Risk factors include urinary tract infections, kidney stones, pregnancy, diabetes, urinary catheters, vesicoureteral reflux in children, and obstruction of the urinary tract. Around a quarter of women who have had one kidney infection have a recurrence within a year.

\section*{Aetiology and Pathophysiology}
Kidney infection usually begins as a bladder infection with bacteria, most commonly Escherichia coli, that ascend the ureters into the kidney. The bacteria multiply in the kidney tissue, causing inflammation, and can enter the bloodstream, producing systemic illness. Factors that promote ascent include obstruction of urine flow, reflux of urine backwards from the bladder, pregnancy, which dilates the ureters, and instruments such as catheters. In men, kidney infection is often linked to prostate enlargement or other obstruction. The inflammation causes the characteristic flank pain and fever, and in severe cases, collections of pus (abscesses) can form, and bacteria in the blood can cause sepsis.

\section*{Clinical Features}
\begin{itemize}
    \item Fever, often high, with chills and sweats
    \item Pain in the flank or lower back, on one or both sides
    \item Tenderness over the kidney
    \item Pain or burning with urination, and urinary frequency
    \item Nausea, vomiting, and feeling generally unwell
    \item Cloudy, smelly, or bloody urine
    \item Headache and fatigue
    \item In older people, confusion may be the main feature
    \item Children may present with fever, vomiting, and poor feeding
\end{itemize}

\section*{Differential Diagnosis}
Other causes of fever and flank pain must be considered.
\begin{itemize}
    \item \textbf{Bladder infection (cystitis):} burning and frequency without high fever or flank pain
    \item \textbf{Kidney stones:} colicky pain, sometimes with fever if infected
    \item \textbf{Kidney abscess:} persistent fever despite antibiotics
    \item \textbf{Bowel conditions:} appendicitis, diverticulitis, and other abdominal infections
    \item \textbf{Back and muscle problems:} musculoskeletal flank pain without fever
    \item \textbf{Gallbladder and pancreatic disease:} upper abdominal pain
    \item \textbf{Sepsis from other sources:} other infections
\end{itemize}

\section*{When to Seek Medical Care}
Anyone with fever and flank pain, or a urinary infection with fever and feeling unwell, should be assessed promptly. Pregnant women, people with diabetes or kidney stones, and older people should seek care early, as they are at higher risk of complications. Urgent care is needed for signs of sepsis, vomiting, or inability to take antibiotics.

\textbf{Contact your local emergency services for:}
\begin{itemize}
    \item High fever with shaking, confusion, or feeling very unwell, which may indicate sepsis
    \item Severe flank or abdominal pain
    \item Vomiting with an inability to keep fluids or medicines down
    \item Difficulty breathing or collapse
    \item Passing very little urine
    \item A kidney infection in pregnancy that is worsening
\end{itemize}

\section*{Diagnosis and Investigations}
\begin{itemize}
    \item \textbf{Urine tests:} dipstick, microscopy, and culture with sensitivity
    \item \textbf{Blood tests:} full blood count, inflammatory markers, kidney function, and blood cultures
    \item \textbf{Imaging:} ultrasound or CT for stones, obstruction, and abscesses
    \item \textbf{Pregnancy test:} in women of reproductive age
    \item \textbf{Assessment of risk factors:} diabetes, reflux, and catheter use
    \item \textbf{Follow-up urine testing:} to confirm clearance
\end{itemize}

\section*{Management}
\begin{itemize}
    \item \textbf{Antibiotics:} oral for mild cases, intravenous for severe cases or in pregnancy
    \item \textbf{Hospital care:} for people who are unwell, vomiting, pregnant, or at high risk
    \item \textbf{Fluids:} intravenous fluids when needed
    \item \textbf{Pain relief:} analgesics and fever control
    \item \textbf{Treatment of complications:} drainage of abscesses and relief of obstruction
    \item \textbf{Management of risk factors:} stones, reflux, and catheters
    \item \textbf{Follow-up:} repeat urine tests and review
    \item \textbf{Prevention of recurrence:} addressing underlying causes
\end{itemize}

\section*{Complications}
\begin{itemize}
    \item Sepsis and septic shock
    \item Kidney abscess and scarring
    \item Permanent kidney damage and impaired function
    \item Recurrent infections
    \item Complications in pregnancy, including preterm birth
    \item High blood pressure from kidney scarring
    \item Rarely, kidney failure with repeated or severe infections
\end{itemize}

\section*{Prognosis and Outcomes}
Most people with kidney infection recover fully with a course of antibiotics, with improvement within a few days and full recovery within weeks. Severe infections, especially in pregnancy, diabetes, or with obstruction, require hospital care but respond well. Complications such as abscess and sepsis are treatable with drainage and intensive care. Managing underlying causes prevents recurrence and long-term kidney damage.

\section*{Prevention and Self-Care}
\begin{itemize}
    \item Treat bladder infections promptly
    \item Drink plenty of fluids
    \item Empty the bladder fully and regularly
    \item Seek medical care early for fever and flank pain
    \item Complete the full course of antibiotics
    \item Attend follow-up urine tests
    \item Manage diabetes and other risk factors
    \item Pregnant women should report urinary symptoms early
    \item Seek urgent care for signs of sepsis
\end{itemize}

\section*{Sources and Further Reading}
The following reputable sources provide further information for healthcare students and patients seeking reliable, current advice about kidney infection.
\begin{itemize}
    \item Kidney Health Australia. \textit{Kidney infection.} https://kidney.org.au/
    \item Healthdirect Australia. \textit{Kidney infection (pyelonephritis).} https://www.healthdirect.gov.au/
    \item Pregnancy, Birth and Baby. \textit{Urinary infections in pregnancy.} https://www.pregnancybirthbaby.org.au/
    \item Better Health Channel. \textit{Kidney infection.} https://www.betterhealth.vic.gov.au/
    \item NHS. \textit{Kidney infection.} https://www.nhs.uk/conditions/kidney-infection/
    \item Mayo Clinic. \textit{Kidney infection.} https://www.mayoclinic.org/
\end{itemize}
